\section{Diskusi}

\subsection{Interpretasi Biologis}

Temuan audit menunjukkan bahwa Shen dan Kukurugya tidak saling
bertentangan --- mereka mengukur besaran yang berbeda. Kukurugya mengukur
\textit{kapasitas maksimal} enzim ($V \cdot \gamma$), sementara Shen
mengukur \textit{efisiensi terealisasi} ($u \cdot V \cdot \gamma$) dengan
$u < 1$ karena ekspresi berlebih konstitutif protein glikolitik.

Fenomena ekspresi berlebih ini bukan artefak pengukuran: Shen melaporkan
bahwa sel mempertahankan protein glikolitik dalam jumlah besar sebagai
cadangan (``\textit{proteome hedging}'') untuk ketahanan terhadap hipoksia.
Implikasinya: pada kondisi fisiologis normal, $u_G$ rendah (banyak enzim
glikolitik idle), sehingga respirasi \textit{terlihat} lebih efisien secara
terealisasi --- persis temuan Shen. Namun, bila diukur pada kapasitas penuh
($u_G = 1$), glikolisis memang lebih cepat per protein --- persis temuan
Kukurugya.

E3b memperkuat temuan ini secara nonlinier: pada keluaran LP (bukan
rumus linier), $u_G$ tetap dominan untuk \textit{E.\ coli} dan ragi,
namun untuk mamalia $\gamma_R$ menjadi parameter terpenting --- menunjukkan
bahwa yield respirasi juga merupakan pengukuran kritis yang selama ini
terselubung oleh analisis linier. Diagram fase E4 menunjukkan bahwa
batas rezim bergeser secara kontinu dengan $u_G$, mengkonfirmasi
bahwa efek \textit{proteome hedging} beroperasi di seluruh rentang
ketersediaan glukosa.

\subsection{Posisi terhadap Wang 2025}

Gambar~\ref{fig:wang} membandingkan kerangka kami dengan Wang secara visual.
Wang memodelkan penyilangan efisiensi sebagai fungsi kualitas nutrisi $q$
--- pada $q$ tinggi respirasi menang, pada $q$ rendah fermentasi menang.
Dalam kerangka kami, variabel yang setara adalah $u_G$: kualitas nutrisi
tinggi $\leftrightarrow$ sel tumbuh lambat $\leftrightarrow$ $u_G$ rendah
(kapasitas glikolitik idle).

Kedua kerangka memprediksi titik penyilangan yang sama tetapi dari sudut
mekanistik yang berbeda. Wang memperkenalkan heterogenitas sel sebagai
mekanisme pendamaian; kami mengidentifikasi $u_G$ sebagai variabel tunggal
yang menggerakkan penyilangan tanpa perlu model heterogenitas. Keduanya
\textbf{komplementer}: Wang menjelaskan \textit{mengapa} $u_G$ bervariasi
antar-sel (heterogenitas), kami menunjukkan bahwa $u_G$ saja cukup untuk
menjelaskan perbedaan verdict.

\begin{figure}[t]
    \centering
    \includegraphics[width=\columnwidth]{F5_wang_comparison.png}
    \caption{Perbandingan kerangka: baris atas = sumbu $u_G$ kami, baris bawah = sumbu kualitas nutrisi $q$ Wang. Titik penyilangan sama, mekanisme berbeda.}
    \label{fig:wang}
\end{figure}

\subsection{Validasi Mekanistik (M6, Peregangan)}

Sebagai validasi tambahan, kami menambahkan kendala massa-enzim pada model
\textit{e\_coli\_core} via COBRApy (Gambar~\ref{fig:eccore}). Ketika
anggaran proteom $\Phi$ diperketat, model FBA menunjukkan onset sekresi
asetat (metabolisme limpahan) --- secara kualitatif konsisten dengan
prediksi model mainan bahwa keterbatasan proteom mendorong transisi
respirasi$\to$fermentasi.

\begin{figure}[t]
    \centering
    \includegraphics[width=\columnwidth]{F4_eccore.png}
    \caption{M6: FBA berbatas-enzim pada \textit{e\_coli\_core}. Kiri: penyapuan $\Phi$ menunjukkan onset limpahan asetat. Kanan: penyapuan glukosa pada $\Phi$ ketat menunjukkan peningkatan fraksi fermentatif.}
    \label{fig:eccore}
\end{figure}

\subsection{Keterbatasan}

Beberapa keterbatasan perlu diakui:
\begin{enumerate}
    \item \textbf{Model \textit{coarse-grained}.} Kedua model adalah
          penyederhanaan dua-sektor; model skala-genom (GECKO, ETFL) akan
          memberikan resolusi lebih tinggi.
    \item \textbf{Parameter provisional.} Nilai parameter ditandai
          PROVISIONAL dan berasal dari ringkasan tabel, bukan pengukuran
          langsung kami. Verifikasi dari SI asli diperlukan.
    \item \textbf{Tiga organisme saja.} Audit terbatas pada
          \textit{E.\ coli}, ragi, dan sel mamalia; generalisasi ke
          organisme lain belum teruji.
    \item \textbf{In-silico.} Seluruh temuan bersifat komputasional;
          validasi eksperimental (pengukuran $u_G$ in-vivo) diperlukan.
\end{enumerate}
